\begin{frame}
    \begin{itemize}
        \item     Soit $d \in \bbN^*$. Soit un ouvert borné $\Omega \subset \bbR^d$.
Pour simplifier l'étude, on prendra des conditions de Dirichlet homogènes et un second membre nul.
Si on se donne une donnée initiale $u_0$, on peut reformuler 
l'équation de Rayleigh-Stokes de la façon suivante : trouver $u \in \scrC^1(\bbR, \Hone)$ 
tel que pour tout $t \in \bbR$
\begin{equation}\label{pb_RS}
\begin{cases}
    u'(t) - \nu\Delta u(t) - \alpha\Delta \big(u'(t)\big) = 0  \\
    u(0) = u_0
\end{cases}
\end{equation}
\item     On définit deux opérateurs $M$ et $L$ par
\begin{eqnarray*}
    M := I - \alpha \Delta \ \ \ \ \ &\text{ et }& \ \ \ \ \ L :=-\nu\Delta
\end{eqnarray*}
de sorte à reformuler le problème \eqref{pb_RS} en : trouver $u \in \scrC^1(\bbR, \Hone)$ tel que pour tout $t \in \bbR$
\begin{equation}\label{RS_forme_operateur}
\begin{cases}
    Mu'(t) + Lu(t)=0  \\
    u(0) = u_0
\end{cases}
\end{equation} 
    \end{itemize}
\end{frame}

\subsection{Extension de $M$ et $L$}

\begin{frame}
\begin{block}{Théorème}
    Soient $(H, \langle \cdot, \cdot \rangle)$ et $(H_0, \langle \cdot, \cdot \rangle_0)$ deux espaces de Hilbert tels que $H$ est un sous-espace dense de $H_0$ et 
\begin{equation}\label{ineq_thm_lax}
    \exists k >0,\ \forall x \in H, \ \norm{x} \geq k \norm{x}_0
\end{equation}
Soit $B$ une forme bilinéaire continue et coercive sur $H$. Alors il existe deux opérateurs linéaires $E$ et $E^*$, dont les domaines de définitions $D$ et $D^*$ sont denses dans $H$, tels que
\begin{enumerate}
    \item[1.] On a
    \begin{eqnarray*}
        \forall x \in D,\ \forall y \in H, \ B(x,y) &=& \ps{Ex}{y}_0 \\
        \forall y \in D^*,\ \forall x \in H, \ B(x,y) &=& \ps{x}{E^* y}_0
    \end{eqnarray*}
    \item[2.] L'image de $E$ est $H_0$, de même pour $E^*$.
    \item[3.] $E$ et $E^*$ admettent des inverses $S$ et $S^*$ qui sont continus de $H_0$ dans $H$.
    \begin{eqnarray*}
        \exists k_1>0,\ \forall z \in H_0, \ \norm{Sz} &\leq& k_1 \norm{z}_0 \\
        \exists k_1^*>0,\ \forall z \in H_0, \ \norm{S^*z} &\leq& k_1^* \norm{z}_0
    \end{eqnarray*}
    Et donc, avec \eqref{ineq_thm_lax}, $S$ et $S^*$ sont continus de $H_0$ dans $H_0$
    %\begin{eqnarray*}
    %    \exists k_2>0,\ \forall z \in H_0, \ \norm{Sz}_0 &\leq& k_2 \norm{z}_0 \\
    %    \exists k_2^*>0,\ \forall z \in H_0, \ \norm{S^*z}_0 &\leq& k_2^* \norm{z}_0
    %\end{eqnarray*}
\end{enumerate}
\end{block}
\end{frame}

\begin{frame}
    \begin{itemize}
    \item Dans notre cas, on considère les espaces de Hilbert $H=H^1_0(\Omega)$ et $H_0=L^2(\Omega)$. On utilisera les inclusions suivantes
\begin{equation*}
    \calD(\Omega) \ \subset \ \Hone \ \subset \ L^2(\Omega)
\end{equation*}
et comme $\overline{\calD(\Omega)}^{L^2(\Omega)}$ = $L^2(\Omega)$, on a bien 
\begin{equation*}
    \overline{\Hone}^{L^2(\Omega)} = L^2(\Omega)
\end{equation*} 
\item De plus, on a 
\begin{equation*}
    \forall \varphi \in \Hone, \ \|\varphi\|_{L^2} \leq \|\varphi\|_{H^1}
\end{equation*}
    \item On va maintenant chercher à établir une formulation faible de notre équation du problème \eqref{pb_RS}. Soit donc $v \in \Hone$. 
\begin{eqnarray*}
    \int_{\Omega}Mu'(t)v + \int_{\Omega} Lu(t)v =0 %&\Longrightarrow&
    %\int_{\Omega} u'(t)v - \alpha \int_{\Omega} \Delta u'(t)v - \nu \int_{\Omega}
    %\Delta u(t)v = 0 \\
    %&\overset{\text{Green}}{\Longrightarrow}& \int_{\Omega} u'(t)v + \alpha \int_{\Omega} \nabla u'(t)\cdot\nabla v + \nu \int_{\Omega}
    %\nabla u(t)\cdot \nabla v = 0 \\
    &\overset{\text{Green}}{\Longrightarrow}& \int_{\Omega} \big(u'(t)v+\alpha\nabla u'(t)\cdot\nabla v\big) \\
    &+& \nu\int_{\Omega} \nabla u(t)\cdot \nabla v = 0
\end{eqnarray*} 
\item On définit ainsi deux formes bilinéaires sur $\Hone$
\begin{eqnarray*}
    b_M(\varphi, \psi) = \int_{\Omega} \big(\varphi\psi+\alpha\nabla \varphi\cdot\nabla \psi\big) \ \ \ \ \ &\text{ et }& \ \ \ \ \
    b_L(\varphi, \psi) = \nu\int_{\Omega} \nabla \varphi\cdot \nabla \psi
\end{eqnarray*} 
\end{itemize}
\end{frame}

\begin{frame}
    \begin{itemize}
\item On obtient ainsi le problème variationnel associé au problème continu \eqref{RS_forme_operateur} : pour tout $t \in \bbR$, trouver $u(t) \in \Hone$ tel que
\begin{equation*}
    \forall v \in \Hone, \ b_M\big(u'(t),v\big) + b_L\big(u(t),v\big) = 0
\end{equation*}
    \item On montre que $b_M$ et $b_L$ sont continues et coercives sur $H^1_0(\Omega)$. 
    \item On peut appliquer le théorème, ainsi il existe deux opérateurs $M_0$ et 
    $L_0$, définis respectivement sur $D(M_0)$ et $D(L_0)$, qui sont denses dans $H^1_0(\Omega)$. 
    \item Le théorème nous assure également de l'existence de $M^{-1}_0$ et 
du fait que cet inverse est continu de $L^2(\Omega)$ dans 
$H^1_0(\Omega)$. 
\item Finalement, \eqref{RS_forme_operateur} se 
réécrit comme le problème suivant
\begin{equation}\label{edo_ML}
\begin{cases}
    u'(t) = -M^{-1}_0 L_0u(t)  \\
    u(0) = u_0
\end{cases}
\end{equation}
\end{itemize}
\end{frame}

\begin{frame}
    \begin{itemize}
        \item Nous voulons désormais étendre $-M^{-1}_0 L_0$ pour que \eqref{RS_forme_operateur} soit bien défini dans $\Hone$. 
        \item L'application du théorème nous a donne deux opérateurs étendus qui sont des bijections
\begin{eqnarray*}
    M_0 : D(M_0) \longrightarrow L^2(\Omega) \ \ \ \ \ &\text{ et }& \ \ \ \ \ L_0 : D(L_0) \longrightarrow L^2(\Omega)
\end{eqnarray*}
Cela justifie que  $-M^{-1}_0L_0 : D(L_0)\longrightarrow D(M_0)$ est une bijection. 
\item On considère l'application $-M^{-1}_0L_0 : \calD(\Omega)
\longrightarrow \Hone$. 
\item On montre que, l'opérateur $-\operateurs$ est continue sur 
$\Hone$. Alors, il existe un unique opérateur 
$B : \Hone \longrightarrow \Hone$ qui prolonge $-\operateurs$. 
    \item Notre équation \eqref{edo_ML} devient donc
\begin{eqnarray}\label{edo_ML_B}
\begin{cases}
    u'(t) = Bu(t)  \\
    u(0) = u_0
\end{cases}
\end{eqnarray}
\end{itemize} 
\end{frame}

\subsection{Existence et unicité}

\begin{frame}
    \begin{block}{Remarque}
        Dans le cas d'un système différentiel linéaire de la forme
\begin{equation*}
    \left\{\begin{array}{c}
        u' + Au = 0 \\
        u(0) = u_0
    \end{array}\right.
\end{equation*}
La solution est de la forme $u = e^{-tA} u_0$. 
Pour résoudre \eqref{edo_ML_B} nous allons définir $\exp(tB)$.
    \end{block}
\begin{itemize}
\item On a $B \in \mathscr{L}(\Hone)$, qui est un espace vectoriel complet. 
\item On peut donc définir l'exponentielle de la façon suivante
\begin{eqnarray*}
\forall t \in \mathbb{R},\  E(t) := \exp(tB) = \sum_{k=0}^{+\infty} \frac{(tB)^k}{k!}
\end{eqnarray*} 
\end{itemize}
    \begin{block}{Remarque}
        \begin{enumerate}
        \item[1.] \label{E_groupe_abelien} $\{ E(t) \ | \ t \in \mathbb{R} \}$ est un groupe abélien. 
        \item[2.] Chaque $E(t)$ est un opérateur bornée de $\Hone$ et la dépendance en $t$ est continue pour la topologie des opérateurs. %En effet, soit 
        %$t_0 \in \bbR$, on a par continuité de $\exp$ et par continuité de $B$
        %\begin{eqnarray*}
        %    \vertiii{E(t)-E(t_0)}_1 &=& \vertiii{\ex^{tB}-\ex^{t_0B}}_1 \\
        %    &=& \vertiii{\ex^{t_0B}\big(\ex^{(t-t_0)B}-I\big)}_1 \\
        %    &\leq& \vertiii{\ex^{t_0B}}_1\vertiii{\ex^{(t-t_0)B}-I}_1 \\
        %    &=& \vertiii{\ex^{t_0B}}_1 \vertiii{\sum_{k=1}^{+\infty}\frac{\big((t-t_0)B\big)^k}{k!}} \\
        %    &\leq& \vertiii{\ex^{t_0B}}_1 (\exp{(|t-t_0|\vertiii{B}_1)}-1)\underset{t \to t_0}{\longrightarrow} 0
        %\end{eqnarray*}
        \item[3.] $E(t)$ est dérivable pour la topologie des opérateurs et $E'(t) = B E(t)$. %En effet, soit $h>0$, on utilise le fait que $\{ E(t) \ | \ t \in \mathbb{R} \}$ est un groupe abélien et que $E$ est continue en $t$
        %\begin{eqnarray*}
        %    \vertiii{\frac{E(t+h)-E(t)}{h}-BE(t)}_1 &=& \vertiii{E(t)\Big(\frac{E(h)-I}{h}\Big)-B} \\
        %    &\leq& \vertiii{E(t)}_1 \vertiii{\frac{E(h)-I}{h}-B}_1 \\
        %    &=& \vertiii{E(t)}_1 \vertiii{BE(h)-B}_1\underset{h \to 0}{\longrightarrow} 0
        %\end{eqnarray*}
        %car par un changement d'indice on a obtenu
        %\begin{equation*}
        %    \frac{E(h)-I}{h} = \frac{1}{h}\sum_{k=1}^{+\infty} \frac{h^kB^k}{k!} = \sum_{k=1}^{+\infty} \frac{h^{k-1}B^k}{k!} = \sum_{k=0}^{+\infty}\frac{h^kB^{k+1}}{k!} = BE(h)
        %\end{equation*}
    \end{enumerate}
    \end{block}
\end{frame}

\begin{frame}{Existence}
    \begin{itemize}
        \item Soit $u_0 \in \Hone$.
        \item On pose pour tout $t \in \bbR,\ u(t) = E(t)u_0$. 
        \item Alors par la remarque précédente, on obtient
\begin{eqnarray*}
    u'(t) = E'(t)u_0 = B E(t)u_0 = Bu(t)
\end{eqnarray*} 
\item Le point $1$ de notre remarque implique aussi que $E(0) = I$, donc
\begin{eqnarray*}
    u(0) = E(0)u_0 = u_0
\end{eqnarray*} 
\item On a bien montré l'existence d'une solution au problème \eqref{edo_ML_B}. 
    \end{itemize}
\end{frame}

\begin{frame}{Unicité}
    \begin{itemize}
        \item Montrons maintenant que cette solution est unique. 
        \item Soit $v(t)$ un autre solution du problème. 
        \item On pose $w$ tel que
\begin{eqnarray*}
    \forall t \in \bbR ,\ w(t) := u(t) - v(t)
\end{eqnarray*} 
\item On vérifie que $w(t)$ est aussi solution du problème \eqref{edo_ML_B} avec pour donnée initiale $w(0)=0$. Posons 
\begin{equation*}
    \forall t \in \bbR, \ \alpha(t) := \|w(t)\|_{H^1}^2=\ps{w(t)}{w(t)}_{H^1}
\end{equation*}
        \item On a $\alpha'(t) = 2\langle Bw(t),w(t)\rangle_{L^2}$. 
        \item En appliquant l'inégalité de Cauchy-Schwarz on aboutit à $\abs{\alpha'(t)} \leq 2\norm{B}_1 \alpha(t)$. 
        \item Ce qui nous permet d'utiliser le lemme de Grönwall
\begin{eqnarray*}
    \forall t \in \bbR, \ \alpha(t) \leq \alpha(0)\exp\bigr(\abs{t}\norm{B}_1\bigl)
\end{eqnarray*} 
\item Or, $w(0) = 0$, donc $\norm{w(0)}^2_{H^1}=0$, ainsi
\begin{eqnarray*}
    \forall t \in \bbR, \ \norm{w(t)}^2_1 = 0 \Longrightarrow w(t) = 0
    \Longrightarrow u(t) = v(t)
\end{eqnarray*}
D'où l'unicité de la solution.
    \end{itemize}
\end{frame}

\subsection{Le cas vectoriel}

\begin{frame}
    \begin{itemize}
        \item On peut également dériver une équation de Rayleigh-Stokes vectorielle, mais dans ce cas, on conserve un terme de pression comme pour le problème de Stokes.
        \begin{equation*}
            \left\{\begin{array}{l}
                \partial_t\bfm{u} - \nu\Delta \bfm{u} - \alpha\Delta \big(\partial_t \bfm{u}\big) = -\nabla \tilde{p} + \bfm{f}  \\
                \Div (\bfm{u}) = 0
            \end{array}\right.
        \end{equation*}
        \item On peut reproduire l'analyse précédente en se plaçant, comme pour le problème de Stokes, sur l'espace à divergence nulle $V := \{\bfm{v} \in \big(\Hone\big)^d \ \ | \ \ \Div(\bfm{v}) = 0\}$.
        \item Il faut cependant utiliser le lemme de De Rham pour reconstruire le champ de pression.
        \begin{block}{Lemme de De Rham}
    Soient un ouvert $\Omega \subset \bbR^d$ et $\bfm{f} \in \big(\calD'(\Omega)\big)^d$. Alors il existe 
    $p \in \calD'(\Omega)$ tel que $\bfm{f} = \nabla p$ si et seulement si 
    \begin{equation*}
        \forall \bfm{v} \in \calD_{\sigma}(\Omega), \ \langle \bfm{f}, \bfm{v} \rangle =0
    \end{equation*}
\end{block}
\item La version la plus générale de cette équation, encore étudiée aujourd'hui, et cette forme vectorielle avec le laplacien fractionnaire.
\begin{equation*}
            \left\{\begin{array}{l}
                \partial_t\bfm{u} + \nu(-\Delta)^{\beta} \bfm{u} + \alpha(-\Delta)^{\beta} \big(\partial_t \bfm{u}\big) = -\nabla \tilde{p} + \bfm{f}  \\
                \Div (\bfm{u}) = 0
            \end{array}\right.
        \end{equation*}
    \end{itemize}
\end{frame}

%\begin{frame}
%\begin{itemize}
%        \item On considère le problème vectoriel : trouver $(\bfm{u},\tilde{p}) \in \scrC^1\big(\bbR, \big(\Hone\big)^d\big) \times \scrC^0(\bbR, L^2(\Omega))$ tel que
%\begin{equation*}
%    \left\{\begin{array}{l}
%        \bfm{u}'(t) - \nu\Delta \bfm{u}(t) - \alpha\Delta \big(\bfm{u}'(t)\big) = -\nabla \tilde{p}(t) \\
%         \Div (\bfm{u}(t)) = 0 \\
%         \bfm{u}(0) = \bfm{u}_0 
%    \end{array}\right.
%\end{equation*}
%\item Pour résoudre ce problème on a cherché à reproduire le 
%raisonnement dans le cas scalaire en adaptant les espaces 
%considérés.
%\item On introduit l'espace de Hilbert
%\begin{equation*}
%    V := \{\bfm{v} \in \big(\Hone\big)^d \ \ | \ \ \Div(\bfm{v}) = 0\}
%\end{equation*}
%et l'espace 
%\begin{equation*}
%    \calD_{\sigma}(\Omega) := \{\bfm{\varphi} \in \big(\calD(\Omega)\big)^d \ \ | \ \ \Div(\bfm{\varphi}) = 0\}
%\end{equation*}
%    \end{itemize}
%\begin{block}{Théorème}
%    Si $\Omega \subset \bbR^d$ est un ouvert borné à la frontière suffisamment régulière alors $\calD_{\sigma}(\Omega)$ est dense dans $V$.
%\end{block}
%\end{frame}

%\begin{frame}{Détermination de la pression}
%    \begin{itemize}
%        \item Pour tout $t \in \bbR$, $\bfm{u}(t)$ est solution au sens des distributions. On a donc pour tout $t \in \bbR$
%\begin{equation*}
%    \forall \bfm{v} \in \calD_{\sigma}(\Omega), \ \langle \bfm{u}'(t) - \nu \Delta \bfm{u}(t) - \alpha \Delta \bfm{u}'(t), \bfm{v} \rangle = 0
%\end{equation*}
%\end{itemize}
%\begin{block}{Lemme de De Rham}
%    Soient un ouvert $\Omega \subset \bbR^d$ et $\bfm{f} \in \big(\calD'(\Omega)\big)^d$. Alors il existe 
%    $p \in \calD'(\Omega)$ tel que $\bfm{f} = \nabla p$ si et seulement si 
%    \begin{equation*}
%        \forall \bfm{v} \in \calD_{\sigma}(\Omega), \ \langle \bfm{f}, \bfm{v} \rangle =0
%    \end{equation*}
%\end{block}
%\begin{itemize}
%\item On applique ce résultat pour $\bfm{u}'(t) - \nu \Delta \bfm{u}(t) - \alpha \Delta\bfm{u}'(t) \in V'$, ce qui nous donne l'existence de $p \in \calD'(\Omega)$ tel que 
%\begin{equation*}
%    \bfm{u}'(t) - \nu \Delta \bfm{u}(t) - \alpha\Delta \bfm{u}%'(t) = -\nabla p
%\end{equation*}
%\item Ceci étant vrai pour tout $t \in \bbR$, on définit une fonction $\tilde{p}:\bbR \rightarrow \calD'(\Omega)$ telle que 
%\begin{equation*}
%    \bfm{u}'(t) - \nu \Delta \bfm{u}(t) - \alpha\Delta \bfm{u}'(t) = -\nabla \tilde{p}(t)
%\end{equation*}
%    \end{itemize}
%\end{frame}